\section{激活函数：为什么需要非线性}

\subsection{从阈值函数到可微激活函数}

神经元模型的历史始于 McCulloch \& Pitts（1943）的阈值单元：

\begin{figure}[H]
\centering
\includegraphics[width=0.95\textwidth]{slides-images/slide-005.jpg}
\caption{生物神经元与 McCulloch-Pitts 阈值单元：$\mathbf{1}(Wx > \theta)$，阈值函数没有斜率，无法进行基于梯度的学习\protect\footnotemark}
\end{figure}
\footnotetext{Slides 第6页。}

阈值函数的问题在于它是分段常数函数——斜率处处为零（除了不连续点），因此\textbf{无法提供梯度信息}来指导学习。梯度学习的核心思想是``找到哪里更陡峭，朝那个方向走''——如同滑雪，你需要坡度来决定方向。

\subsection{常见激活函数一览}

\begin{figure}[H]
\centering
\includegraphics[width=0.95\textwidth]{slides-images/slide-006.jpg}
\caption{激活函数家族：从经典的 logistic/tanh 到现代的 ReLU/Swish/GELU\protect\footnotemark}
\end{figure}
\footnotetext{Slides 第7页。}

\textbf{Logistic（Sigmoid）函数}：
$$\sigma(z) = \frac{1}{1 + \exp(-z)}$$
输出范围 $(0, 1)$，常用于输出层（映射到概率）。缺点：输出始终为正，且在两端梯度趋近于零（梯度消失）。

\textbf{Tanh 函数}：
$$\tanh(z) = \frac{e^z - e^{-z}}{e^z + e^{-z}}$$
输出范围 $(-1, 1)$，是 sigmoid 的缩放平移版本：$\tanh(z) = 2\sigma(2z) - 1$。

\textbf{ReLU（Rectified Linear Unit）}：
$$\text{ReLU}(z) = \max(z, 0)$$
正半轴斜率恒为 1，计算极快。虽然负半轴``死区''（梯度为零）看似矛盾，但实践中效果极好——它提供了一种自然的稀疏激活和特化机制。

\begin{warningbox}{ReLU 的``死区''问题}
当 ReLU 神经元的输入持续为负时，梯度为零，该神经元``永久死亡''，无法再被激活。这就是所谓的 \textbf{dying ReLU} 问题。虽然在正常训练中这通常不是致命问题（网络中总有一部分神经元是活跃的），但在学习率过大或初始化不当时可能导致大量神经元死亡。
\end{warningbox}

\textbf{Leaky ReLU / Parametric ReLU}：在负半轴给予一个小斜率（如 0.01），避免完全死亡。Parametric ReLU 将负半轴斜率作为可学习参数。

\textbf{Swish 和 GELU}：近年来在 Transformer 模型中广泛使用。
$$\text{Swish}(x) = x \cdot \sigma(x), \qquad \text{GELU}(x) = x \cdot P(X \le x), \; X \sim \mathcal{N}(0,1)$$

两者形状类似：正半轴近似 $y = x$，负半轴有一小段``弯曲''区域提供梯度。

\subsection{为什么非线性不可或缺}

\begin{figure}[H]
\centering
\includegraphics[width=0.95\textwidth]{slides-images/slide-007.jpg}
\caption{非线性的必要性：没有非线性，多层矩阵乘法 $W_1 W_2 x = Wx$ 只是单个线性变换；有了非线性，网络可以逼近任意复杂函数\protect\footnotemark}
\end{figure}
\footnotetext{Slides 第8页。}

\begin{importantbox}{线性层的叠加 = 单层线性变换}
如果只做矩阵乘法而没有非线性：
$$h_2 = W_2(W_1 x) = (W_2 W_1) x = W'x$$
多层网络在表示能力上与单层完全等价！非线性激活函数是使神经网络成为\textbf{万能函数逼近器}（universal function approximator）的关键。
\end{importantbox}

\begin{knowledgebox}{线性网络在学习理论中的价值}
Manning 教授指出了一个有趣的现象：虽然线性神经网络在\textbf{表示能力}上没有优势（等价于单层线性变换），但在\textbf{学习动力学}上却有独特性质。理论界有相当多论文研究线性网络的学习行为，因为多层结构即使没有非线性也会影响优化路径。
\end{knowledgebox}

\subsection{本章小结}

激活函数是神经网络能够学习复杂非线性映射的关键。从历史上的阈值函数到 sigmoid/tanh，再到 ReLU 及其变体，直到 Transformer 时代的 Swish/GELU，激活函数的演进反映了对``梯度流通性''和``计算效率''的不断追求。核心原则是：保留足够的梯度信号使学习成为可能。

%% ========== Part 4: 随机梯度下降回顾 ==========

